\section{Tinjauan Pencapaian Kompetensi Secara Menyeluruh}

\subsection{Pemetaan Asesmen ke CPMK}

\begin{table}[H]
\centering
\begin{tabular}{|>{\raggedright\arraybackslash}p{5cm}|>{\raggedright\arraybackslash}p{8cm}|}
\hline
\textbf{CPMK / Sub-CPMK} & \textbf{Diukur Melalui} \\
\hline
CPMK-1: Sub-CPMK 1.1--1.3 (Proposisi, operator, tabel kebenaran, tautologi, inferensi) & Pilihan Ganda (1--3, 5), Tugas Praktik (1), Bagian C (1) \\
\hline
CPMK-2: Sub-CPMK 2.1--2.3 (Logika predikat, kuantor, translasi, pemodelan) & Pilihan Ganda (4), Tugas Praktik (2, 5), Bagian D \\
\hline
CPMK-3: Sub-CPMK 3.1--3.3 (Aljabar Boolean, K-Map, gerbang logika) & Pilihan Ganda (8), Tugas Individu 2, Bagian C (jika ada soal K-Map) \\
\hline
CPMK-4: Sub-CPMK 4.1, 4.2, 4.3 (Metode bukti, induksi, verifikasi program) & Pilihan Ganda (6, 7, 9), Tugas Praktik (3--6), Bagian C (2--3), Bagian D \\
\hline
\end{tabular}
\caption{Pemetaan Komponen Asesmen ke CPMK dan Sub-CPMK}
\end{table}

\subsection{Self-Assessment Checklist}

Sebelum mengerjakan asesmen akhir, verifikasi penguasaan atas hal-hal berikut:

\begin{checklist}
  \item Proposisi dan operator logika (negasi, konjungsi, disjungsi, implikasi, bikondisional)
  \item Tabel kebenaran dan ekuivalensi logika
  \item Tautologi, kontradiksi, dan kontingensi
  \item Translasi bahasa alami ke simbol logika proposisional dan predikat
  \item Kuantifier universal dan eksistensial beserta negasinya
  \item Modus ponens, modus tollens, silogisme, dan aturan inferensi lain
  \item Bukti langsung, kontraposisi, dan kontradiksi
  \item Induksi matematika (basis dan langkah induktif)
  \item Representasi ekspresi Boolean ke gerbang logika dasar (AND, OR, NOT).
  \item Pemodelan himpunan fakta, aturan generalisasi, dan penarikan kesimpulan basis pengetahuan.
  \item Verifikasi kebenaran program: prekondisi, postkondisi, loop invariant, kebenaran parsial dan total, serta pembuktian terminasi (Sub-CPMK 4.3).
\end{checklist}
